\documentclass{article}
\usepackage[T1]{fontenc}
\usepackage{lmodern}
\usepackage{graphicx}
\usepackage{amsmath,amssymb}
\usepackage[a4paper,margin=2cm]{geometry}
\usepackage[absolute,overlay]{textpos}
\setlength{\TPHorizModule}{1mm}
\setlength{\TPVertModule}{1mm}
\usepackage[indonesian]{babel}
\usepackage{hyperref}
\setlength{\emergencystretch}{2em}
% unit: Halaman 12

\begin{document}

% === Halaman 12 (python/native) ===

12

Prosedur Dekripsi

1.

Gunakan kunci privat x untuk menghitung (ax)–1 = ap – 1 – x mod p 

2.

Hitung plainteks m dengan persamaan:


 m = b/ax mod p = b(ax)–1 mod p

\end{document}
